\documentclass[11pt]{article}
\usepackage[a4paper,margin=2.4cm]{geometry}
\usepackage{amsmath,amssymb}
\usepackage{enumitem}
\usepackage{hyperref}

\newcommand{\lam}[1]{\lambda #1.\,}
\newcommand{\FV}{\mathrm{FV}}
\newcommand{\reduce}{\twoheadrightarrow_\beta}
\newcommand{\onereduce}{\to_\beta}
\newcommand{\subst}[2]{[#1/#2]}

\title{Lecture 1 In-Class Worksheet\\Untyped Lambda Calculus}
\author{}
\date{}

\begin{document}
\maketitle

\section*{Learning objectives}
By the end of this three-hour unit, students should be able to:
\begin{enumerate}[nosep]
  \item parse lambda terms using the binding and application conventions;
  \item identify free and bound variables;
  \item perform simple capture-avoiding substitutions;
  \item reduce a term by one or more beta steps;
  \item distinguish $=_\alpha$, $\onereduce$, $\reduce$, and $=_\beta$.
\end{enumerate}

\section*{Suggested pacing}
\begin{center}
\begin{tabular}{r p{0.68\linewidth}}
0--20 min & Motivation, syntax, and parsing conventions.\\
20--45 min & Abstract syntax trees and binding.\\
45--65 min & Exercise 1: parsing and binding.\\
65--75 min & Break.\\
75--105 min & Free variables and capture-avoiding substitution.\\
105--135 min & Exercise 2: substitution.\\
135--165 min & Beta-reduction and multi-step reduction.\\
165--180 min & Exit check and summary of relations.\\
\end{tabular}
\end{center}

\section*{Exercise 1: parsing and binding}
Expand the conventions and mark the binders for each occurrence of a variable.
\begin{enumerate}
  \item $\lam{x y} x\;y\;z$
  \item $(\lam{x}x)\;(\lam{y}y)$
  \item $\lam{a b} a\;(\lam{c}a\;b)$
\end{enumerate}

\paragraph{Solution checkpoints.}
\begin{enumerate}[nosep]
  \item $\lam{x}(\lam{y}((x\;y)\;z))$; $z$ is free.
  \item Application of two abstractions; both $x$ and $y$ are bound in their own bodies.
  \item $\lam{a}(\lam{b}(a\;(\lam{c}(a\;b))))$; $c$ does not occur in the body.
\end{enumerate}

\section*{Exercise 2: free variables and substitution}
Compute the free variables first, then perform the substitution if it is safe.
\begin{enumerate}
  \item $\FV(\lam{x}(x\;z))$
  \item $(\lam{y}(x\;y))\subst{z}{x}$
  \item $(\lam{z}(x\;z))\subst{z}{x}$
  \item $(\lam{y}(x\;y))\subst{y}{x}$
\end{enumerate}

\paragraph{Solution checkpoints.}
\begin{enumerate}[nosep]
  \item $\{z\}$.
  \item $\lam{y}(z\;y)$.
  \item The binder $z$ conflicts with the substituting term; alpha-rename first, e.g. $\lam{w}(x\;w)$, then substitute to obtain $\lam{w}(z\;w)$.
  \item Alpha-rename $y$ first, e.g. $\lam{w}(x\;w)$, then substitute to obtain $\lam{w}(y\;w)$.
\end{enumerate}

\section*{Exercise 3: beta-reduction}
Circle every beta-redex and reduce to normal form when possible.
\begin{enumerate}
  \item $(\lam{x}x)\;z$
  \item $((\lam{x}x)\;y)\;((\lam{z}z)\;x)$
  \item $(\lam{x y}x)\;u\;v$
  \item $(\lam{x y}y)\;x$
\end{enumerate}

\paragraph{Solution checkpoints.}
\begin{enumerate}[nosep]
  \item $z$.
  \item $y\;x$ after reducing both visible redexes.
  \item $u$.
  \item $\lam{y}y$ after avoiding capture by alpha-renaming if necessary.
\end{enumerate}

\section*{Exit check}
Students should be able to answer these without notes:
\begin{enumerate}[nosep]
  \item Give two alpha-equivalent but syntactically different terms.
  \item Give a substitution example where alpha-renaming is necessary.
  \item Explain why $\onereduce$ is directed but $=_\beta$ is symmetric.
\end{enumerate}

\section*{Instructor notes}
For a no-homework version of the course, treat Church encodings, the Y combinator, and confluence as optional or appendix material unless the class finishes the core exercises early.

\end{document}
